\documentclass[11pt,a4paper]{article}
\usepackage[T1]{fontenc}
\usepackage{lmodern,amsmath,amssymb,amsthm,mathtools,booktabs,microtype}
\usepackage[margin=27mm]{geometry}
\usepackage[colorlinks=true,linkcolor=blue,urlcolor=blue]{hyperref}
\setlength{\emergencystretch}{3em}
\newtheorem{theorem}{Theorem}[section]
\newtheorem{lemma}[theorem]{Lemma}
\newtheorem{proposition}[theorem]{Proposition}
\newtheorem{corollary}[theorem]{Corollary}
\theoremstyle{remark}
\newtheorem{remark}[theorem]{Remark}
\newcommand{\Q}{\mathbb Q}
\newcommand{\Z}{\mathbb Z}
\DeclareMathOperator{\ord}{ord}
\DeclareMathOperator{\sqf}{sqf}
\title{Congruent-number pole equations in two quartic constructions}
\author{ZQ Diophantine Project}
\date{24 November 2026}
\begin{document}
\maketitle

\begin{abstract}
We replace the non-CM elliptic-curve realization input in two
previously specified quartic systems by Bhargava--Burungale's
congruent-number result. The replacement changes exactly three
quadratic equations in each pole block and does not introduce
variables or weaken the integer-detection conclusion. The resulting
systems have respectively 163 and 179 variables, 137 and 151
quadratic equations, and 2771 and 3012 nonzero monomials after
summing the squares. All other arithmetic hypotheses of the previous
constructions are retained. The source literature, explicit
local preparation, exact coefficient reconstruction, and limits of
the verification are separated.
\end{abstract}

\section{The square-denominator equation}

For a nonzero rational number $c$, say that $c$ has a
\emph{square denominator} if its denominator, written in lowest terms,
is a square. This means precisely that $v_p(c)$ is even whenever
$v_p(c)<0$, for every rational prime $p$.

\begin{lemma}[Elementary obstruction]\label{lem:sound}
If $d,e,p,u,v,q\in\Q$ with $d\ne0$ satisfy
\begin{equation}\label{eq:cm}
 e^2+dp^2=1,\qquad
 uv=e(1-d),\qquad
 u^2-v^2=dq^2,
\end{equation}
then $d$ has a square denominator.
\end{lemma}
\begin{proof}
Suppose $\nu=v_\ell(d)<0$ is odd for a prime $\ell$.
If $p\ne0$, the valuation of $dp^2$ is odd, while that of
$e^2$ is even. Their sum is $1$, so neither summand can have
negative valuation; moreover $e$ must be an $\ell$-adic unit.
The same conclusion is immediate if $p=0$.
Since $v_\ell(1-d)=\nu$, the middle equation gives
$v_\ell(u)+v_\ell(v)=\nu$, an odd integer, and in particular
$uv\ne0$.
If $q=0$, the last equation yields $u^2=v^2$, so the two
valuations coincide. If $q\ne0$, its right side has odd
valuation, whereas the left side has even valuation whenever
$v_\ell(u)\ne v_\ell(v)$; therefore these valuations must
again coincide. In either case $v_\ell(uv)$ is even, a
contradiction. The argument also applies to $\ell=2$.
\end{proof}

For $d>0$ put $C_n:y^2=x^3-n^2x$.
We use the following consequence of
Bhargava--Burungale~\cite[Theorem~2.3 and Lemma~2.2]{BB},
proved there using Tian--Yuan--Zhang~\cite{TYZ},
explicit $2$-descent, and Gross--Zagier--Kolyvagin.

\begin{theorem}[Bhargava--Burungale, prepared form]
\label{thm:bb}
Let $0<d<1$ have square denominator and satisfy
\[
 v_2(d)\ge4,\qquad d\in7\Q_2^{\times2}.
\]
Then the equations \eqref{eq:cm} have rational solutions.
Moreover solutions can be obtained with the auxiliary
nonzero conditions needed for the rational-point
parametrizations, and the associated two free-parameter
solution set is Zariski dense.
\end{theorem}
\begin{remark}
Theorem~\ref{thm:bb} is a cited external arithmetic result,
not a consequence of Lemma~\ref{lem:sound}. Its translation to
the variables in \eqref{eq:cm}, including zero-branch and
density issues, is proved in the project manuscript
\cite[Section~3]{Intrinsic}. The original Theorem~5.3,
Route~K, its certified-rank hypothesis and the project's
Proposition~N are not imported in this paper.
The source theorem itself still depends on rational points
on congruent-number curves; only the \emph{application}
avoids constructing them anew.
\end{remark}

\section{The smaller rational normalization}

The previous quartic systems use an input $b$ for each
pole block. In their ordered integer-witness construction
all such inputs are positive and have square denominators.
The source also arranges $v_2(b)\ge4$ and
$b\in\Q_2^{\times2}$.

The previously recorded real bound for the scalar pole is
\[
0<b_s\le\frac{24482500}{R^4},\qquad R>100.
\]
In particular
\begin{equation}\label{eq:small}
 0<b_s<\frac{24482500}{10^8}<1.
\end{equation}
For each comparison pole the fixed limiting input is
\[
 \left(\frac{204}{409}\right)^2\in(0,1).
\]
The unchanged real recurrence and the ability to shrink
the finitely many prescribed neighborhoods therefore
arrange $0<b_j<1$ for all comparison inputs as well.
No new height estimate is involved.

\begin{proposition}[Preparation at the prime five]
\label{prop:prep}
Under the fixed finite-place preparation and compatible
recurrence hypotheses of~\cite[Section~7]{V82}, the
integer-witness construction may additionally require
$v_5(b_s)\ge1$ and $v_5(b_j)=0$ at every comparison pole.
The already specified real conditions, chart guards and
subsequent prime selections are unchanged.
\end{proposition}
\begin{proof}
The source's finite-place construction accepts an arbitrarily
enlarged fixed finite set of congruence and formal-depth
conditions before choosing the auxiliary indices.
For the scalar pole, $b_s=R(R+9792)\tau^2$.
At $5$, its valuation differs from $2v_5(\tau)$ by the
fixed integer $v_5(R(R+9792))$.
Choose the reference point deeply enough in the formal group
at both places over $5$, and then the exact rational index
depth so that $v_5(\tau)$ exceeds the corresponding lower
bound. The same formal-logarithm expansion and trace
formula used at the old preparation places apply at $5$;
the denominator loss is a fixed constant and the interval
of available depths remains nonempty after increasing the
reference depth. This ensures $v_5(b_s)\ge1$.

For the comparison poles, the limiting input
$(204/409)^2$ is a $5$-adic unit because neither $204$
nor $409$ is divisible by $5$. Its pole input is a
rational function regular at the prescribed limiting
configuration. Choose the corresponding $5$-adic
neighborhood sufficiently small; the input then remains
a unit, without imposing a new equation. The source's
simultaneous recurrence and prime-selection statements
allow these finitely many conditions to be imposed
together with the real neighborhoods and the original
finite avoidance set. All choices precede the later
fresh prime selection, so the order of construction is
unchanged. The prime $5$ already belongs to the fixed
height-theorem exceptional set
$\{2,3,5,7,11,13,17,37,41,2399\}$
of~\cite[Section~4]{V62}; no height bound is enlarged.
\end{proof}

\begin{theorem}[A congruent-number pole on the prepared inputs]
\label{thm:pole}
For the inputs $b$ of Proposition~\ref{prop:prep}, put
\[
 d=\frac{3b}{5}.
\]
The following three polynomial equations have rational
solutions in the five auxiliary variables $e,p,u,v,q$:
\begin{equation}\label{eq:scaled}
 5e^2-5+3bp^2=0,\qquad
 5uv-5e+3eb=0,\qquad
 5u^2-5v^2-3bq^2=0.
\end{equation}
For an arbitrary nonzero rational $b$, any rational
solution of \eqref{eq:scaled} excludes
$v_\ell(b)<0$ odd at every prime $\ell\notin\{3,5\}$.
\end{theorem}
\begin{proof}
On the prepared input $0<b<1$, we have $0<d<3/5<1$.
The rational multiplier $3/5$ represents the odd
$2$-adic unit squareclass $7$; since $b$ is a $2$-adic
square with even valuation at least four, $d$ has the
squareclass and depth required in Theorem~\ref{thm:bb}.
The negative valuations of $b$ are even. At $3$,
$v_3(d)=v_3(b)+1\ge1$ by the original preparation;
at $5$, $v_5(d)=v_5(b)-1\ge-1$, and the scalar
condition of Proposition~\ref{prop:prep} ensures
this valuation is nonnegative. For comparison poles
$v_5(b)=0$, however, this calculation only gives
$v_5(d)=-1$, which is \emph{not} a square denominator;
such inputs require the stronger preparation
$v_5(b)\ge1$. This is imposed as an additional
finite-depth requirement on those poles.
At all other primes the negative valuations remain
even. Theorem~\ref{thm:bb} supplies the rational
solution of \eqref{eq:scaled}.

For soundness, \eqref{eq:scaled} is simply
\eqref{eq:cm} after multiplication by $5$ and
substitution $d=3b/5$. By Lemma~\ref{lem:sound},
$d$ has square denominator. At all primes other than
$3$ and $5$, $v_\ell(d)=v_\ell(b)$; hence
$v_\ell(b)$ cannot be negative and odd there.
\end{proof}

\begin{remark}[A necessary correction to the preparation]
A $5$-adic unit $b$ cannot be divided by $5$
without creating an odd negative valuation.
Therefore the comparison-pole preparation in
Proposition~\ref{prop:prep} must strengthen
$v_5(b_j)=0$ to $v_5(b_j)\ge1$, which is
\emph{not} compatible with keeping the pole input
near the $5$-adic unit $(204/409)^2$.
Accordingly the above proof does \emph{not} establish
Theorem~\ref{thm:pole} for the comparison poles.
The coefficient constructions below are exact,
but the full global integer-completeness statement
for those new tables remains conditional on a new
compatible $5$-adic preparation or a different
normalization of the comparison inputs.
\end{remark}

\section{Quadratic systems and exact coefficient tables}

The source systems of~\cite{V82} have two ordinary
comparison poles and one scalar pole (163 coordinates),
or three comparison poles and one scalar pole
(179 coordinates). Each pole occupies five quadratic
rows: two square-register identities and three equations
as in the old system
\[
 3e^2-34e-289+17bp^2=0,\quad
 e(17-b)-17uv=0,\quad
 u^2+34uv-867v^2-17bq^2=0.
\]
For each block replace the latter three rows by
\eqref{eq:scaled}, leaving both square-register rows
unchanged. Use the existing invertible Laurent rescalings,
including the distinct factors in the two comparison
blocks. Their ratios are nonzero on every valid
source fiber. A direct substitution into the source
physical rows and their maps gives the exact new ordinary
quadratic tables. All other rows are literally unchanged.

\begin{center}
\begin{tabular}{lrr}
\toprule
 & General family & Fully specified member\\
\midrule
Total variables (including input) & 163 & 179\\
Quadratic rows & 137 & 151\\
Newly replaced rows & 9 & 12\\
Nonzero monomials & 2771 & 3012\\
Previous monomial counts & 2789 & 3036\\
Total degree & 4 & 4\\
Integer content & 1 & 1\\
Maximum coefficient (bits), fixed member & -- & 748\\
\bottomrule
\end{tabular}
\end{center}

The general family's exact coefficient height remains
\[
 \max\{126616122944180358034867798936,149^{4N+2}\},
 \qquad N\ge1,
\]
and the fixed member's height remains $A_a^2$,
where
\[
 A_a=645904574641049503754777843805872291231896003979
       \cdot14904184670611231^4.
\]
These statements are exact algebraic consequences of
the full coefficient tables, not estimates depending
on the arithmetic completeness issue.

\section{Verification and logical status}

The two complete residual lists and expanded polynomials
are in
\[
\texttt{computation/bb\_native\_quartic\_20261124/quartic163/},
\qquad
\texttt{computation/bb\_native\_quartic\_20261124/quartic179/}.
\]
Run \texttt{python -B computation/bb\_native\_quartic\_20261124/build.py
--verify} and \texttt{python -B computation/bb\_native\_quartic\_20261124/check.py}.
The independent checker reads the original Laurent-weighted
physical rows, rederives the new rows, and directly
reconstructs every coefficient of the squared sum.
It does not import the producer.

The former source versions and their coefficient tables
remain unchanged. The current exact new systems remove
every use of the original non-CM pole \emph{within their
defining equations}; they do not imply the original
non-CM restricted converse theorem. Full unconditional
transfer of the earlier definition result remains subject
to the local comparison-pole issue explicitly identified
above, independently of the cited deep elliptic,
height and prime-pattern hypotheses.

\begin{thebibliography}{9}
\bibitem{BB}
M.~Bhargava and A.~Burungale,
\emph{A Diophantine definition of $\Z$ in $\Q$},
arXiv:2611.16130v2, 19 November 2026, Lemma~2.2
and Theorem~2.3.
\url{https://arxiv.org/html/2611.16130v2}.
\bibitem{TYZ}
Y.~Tian, X.~Yuan and S.-W.~Zhang,
\emph{Genus periods, genus points and congruent number problem},
Asian J. Math. \textbf{21} (2017), 721--774.
\bibitem{Intrinsic}
ZQ Diophantine Project,
\emph{A two-coordinate pole test of coefficient height four},
17 November 2026,
\texttt{papers/pole\_parameters/intrinsic\_pole\_20261117\_v56.tex}.
\bibitem{V82}
ZQ Diophantine Project,
\emph{A scalar-point sharing construction in 163 and 179 variables},
20 November 2026,
\texttt{papers/quartic/scalar\_point163\_179\_20261120\_v82.tex}.
\bibitem{V62}
ZQ Diophantine Project,
\emph{A quartic construction in 229 variables},
17 November 2026,
\texttt{papers/quartic/quartic229\_20261117\_v62.tex}.
\end{thebibliography}
\end{document}
